\section{Related Work}
\label{sec:related}

\para{Catastrophic forgetting and continual learning.}
Sequential training of neural networks overwrites earlier knowledge~\cite{mccloskey1989catastrophic}.
A large body of work proposes remedies, such as regularizing important weights~\cite{kirkpatrick2017overcoming} or constraining updates with stored examples~\cite{lopezpaz2017gradient}, and summarizes outcomes with average accuracy and backward transfer~\cite{lopezpaz2017gradient}.
We study the unmitigated case, naive sequential fine-tuning, because it is the standard lower bound and because our question is how forgetting unfolds, not how to prevent it.

\para{Continual instruction tuning of LLMs.}
Recent benchmarks extend continual learning to instruction-tuned language models~\cite{scialom2022fine,zhang2023citb,wang2023trace,luo2023empirical}.
These studies establish that forgetting is substantial in this setting and compare mitigation methods across task streams.
They mostly report end-of-stream aggregates under a small number of fixed task orders.
Unlike them, we evaluate at sub-stage resolution, vary the order with a balanced design~\cite{williams1949experimental}, and fit functional forms to the trajectories.
Our use of LoRA~\cite{hu2022lora} follows common practice; \citet{biderman2024lora} compare forgetting under LoRA and full fine-tuning, a comparison that our study does not include.

\para{The time course of forgetting.}
\citet{delange2023continual} evaluate continually during training and report a transient but large drop on old tasks right after a task switch, which end-of-task evaluation misses.
Our onset experiment observes a fast drop of the same kind in LoRA instruction tuning, here without recovery inside the stage.
\citet{ramasesh2021anatomy} relate forgetting to task similarity, which is consistent with the pair-specific structure we find.
\citet{kotha2024understanding} argue that part of what fine-tuning ``forgets'' is the model's inference of which task it is being asked to perform; our observation that old tasks lose their output \emph{format} after later training fits this reading.

\para{Power laws of forgetting.}
In human memory research, retention over time is often described by a power function~\cite{wixted1991form}.
\citet{anderson1997artifactual} show that averaging non-power curves across items or subjects can produce a power-like aggregate.
We find the same pattern in our data: the power law fits averaged curves on one model while saturating curves fit the individual trajectories.
We did not test the aggregation account directly.
